\section{A staircase made by sampling the digits}
\label{samp:section}

Reading the digits at equally spaced positions gives another real number:
\[
 Y_g(a,d)=\sum_{j\ge0}x_g(a+jd)2^{-j-1}.
\]
Now choose the starting point $a$ and the step $d$ uniformly from
$1,\ldots,N$, subject to $\gcd(a,d)=1$.  This last restriction prevents
the entire progression from being trapped in a common prime lane.
Let $\nu_{g,N}$ be the distribution of $Y_g(a,d)$.  Its limiting
distribution gives a curve that rises across every interval, even in
examples where its derivative is zero almost everywhere.  The apparent
contradiction disappears once we distinguish total rise from the slope
at a typical point.

Write $\mathcal F=\{p_j:g_j=0\}$ for the disabled primes.  The hypothesis
needed below says that the successful primes retain a divergent
reciprocal sum:
\begin{equation}
 \sum_{p\notin\mathcal F}\frac1p=\infty.
 \label{samp:success-hypothesis}
\end{equation}
It holds unconditionally for the actual Goldbach mask.  Indeed, let
$E(Y)$ count even Goldbach exceptions in $[4,Y]$ and let
$A(R)=\#(\mathcal F\cap[2,R])$.  Pintz's Theorem~1
\cite{sampPintz2018} states that $E(Y)<Y^{18/25}$ beyond an
ineffective threshold.  The prime number theorem therefore gives
\begin{equation}
 A(R)=E\bigl(2(\pi(R)+1)\bigr)
 \ll (R/\log R)^{18/25} \quad(R\ge3).
 \label{samp:pintz}
\end{equation}
Partial summation shows $\sum_{p\in\mathcal F}1/p<\infty$, whereas
the sum over all primes diverges.  We use this established estimate
as an input.  Its precise exponent is useful later; here any power
strictly smaller than one would suffice.

\begin{theorem}[Primitive sampling]\label{samp:staircase}
Assume \eqref{samp:success-hypothesis}.  Then $\nu_{g,N}$ converges
weakly to a nonatomic probability measure $\nu_g$ of full support
$[0,1]$ and Hausdorff dimension one.  Its Lebesgue decomposition is
\begin{equation}
 \nu_g=L_g\,\mathrm{Leb}+\nu_{g,\mathrm{sc}},\qquad
 L_g=\prod_{p\in\mathcal F}\frac1{p+1},
 \label{samp:decomposition}
\end{equation}
where $\nu_{g,\mathrm{sc}}$ is singular continuous.  The empty product
is one; an infinite product here is zero.  Consequently
\[
 G_g(y):=\nu_g([0,y])\quad(0\le y\le1)
\]
is a homeomorphism of $[0,1]$ onto itself and
$G_g'(y)=L_g$ for Lebesgue-almost every $y$.
\end{theorem}

Full support means that every nonempty open interval receives positive
probability.  Singular means that a set of ordinary length zero carries
the measure.  Dimension one means that every set carrying its full mass
has Hausdorff dimension one.  These properties can coexist: small in
length need not mean small in dimension.  We do not claim that the
measure has the same local dimension at almost every point.

The mechanism is easier to see than the full convergence proof.
Choose, independently at every prime $p$, a residue pair $(U,V)\bmod p$
uniformly among the $p^2-1$ pairs other than $(0,0)$.  These simultaneous
residues are a primitive profinite pair.  Independently choose fair
bits $B_0,B_1,\ldots$.  The limiting sampled real has the law of
\[
 Z_g=\sum_{j\ge0}\bigl(1-h_g(U+jV)\bigr)B_j2^{-j-1}.
\]
Thus an open lane contributes a fair bit and a disabled lane contributes
zero.  This is a theorem about the average over starting points and
steps; it does not assume that the original consecutive digits are
independent.

A lane $p$ is visited precisely when $p$ does not divide $V$.  The
probability that it is avoided is $1/(p+1)$.  Consequently every bit is
free precisely on an event of probability $L_g$; that event supplies the
uniform part of the measure.  Once a disabled lane is visited, zeros
repeat on an infinite arithmetic progression.  Such digit strings form
a single Lebesgue-null set, which carries the remaining part.

This explains the decomposition and its product formula.  Full support
comes from making any finite prefix open by finitely many compatible
congruences.  Full dimension requires an additional limiting argument:
there are positive portions of the law with a proportion of free bits
arbitrarily close to one.  The complete proof, including the
Frantzikinakis--Host variable-step cancellation that produces the fair
bits, is in Appendix~\ref{samp:proofs}.


\begin{corollary}\label{samp:goldbach-trichotomy}
For the original Goldbach constant, $G_{\Gold}$ is the identity if and
only if Goldbach holds.  If there are finitely many exceptions and at least one, then
\[
 G_{\Gold}'(y)=\frac1{\prod_{p\in\mathcal F}(p+1)}\in(0,1)
 \quad\hbox{almost everywhere}.
\]
There are infinitely many Goldbach exceptions if and only if
$G_{\Gold}'=0$ almost everywhere.  In this last case $G_{\Gold}$ is a singular
homeomorphism whose associated measure has dimension one.
\end{corollary}

This phenomenon also has an unconditional explicit example.  Disable
exactly the lanes
\begin{equation}
 \mathcal F_* =\{p_{2^j}:j\ge2\}.
 \label{samp:explicit-mask}
\end{equation}
This is a concrete mask, defined without asking any unresolved arithmetic
question.  Its disabled-prime reciprocal sum converges because $p_{2^j}>2^j$.  The associated
$G_*$ is therefore a strictly increasing singular homeomorphism,
with full-support measure of dimension one.  Thus the geometry
exists independently of which Goldbach alternatives occur.
